\chapter{Spark SQL and Delta Lake Quick Reference}
\index{Spark SQL}\index{Delta Lake!reference}\index{PySpark!reference}%
This appendix consolidates the Spark, SQL, and Delta Lake syntax used across the
book into one lookup. Examples run on Azure Databricks or Synapse Spark unless
noted; T-SQL fragments are Synapse dedicated or serverless pool syntax.

\section{Reading and Writing}
\index{Spark SQL!read/write}

\begin{lstlisting}[language=Python,caption={The lake read/write patterns used throughout Parts II--V.},label={lst:app-spark-io}]
LAKE = "abfss://curated@<adls>.dfs.core.windows.net"

df = spark.read.parquet(f"{LAKE}/orders/")
df_json = (spark.read.schema("order_id STRING, amount DOUBLE")
           .json(f"{LAKE}/raw/orders/"))
df_csv = (spark.read.option("header", True)
          .option("inferSchema", False)      # declare schemas; inference costs a pass
          .csv(f"{LAKE}/raw/orders/"))

(df.write.format("delta").mode("overwrite")
   .option("overwriteSchema", "true")
   .partitionBy("order_year", "order_month")
   .save(f"{LAKE}/orders_v2/"))

spark.sql(f"CREATE TABLE IF NOT EXISTS orders_v2 "
          f"USING DELTA LOCATION '{LAKE}/orders_v2/'")
\end{lstlisting}

\section{DataFrame Operations and the Shuffle Boundary}
\index{shuffle!reference}

\begin{center}
\footnotesize
\begin{tabular}{@{}p{5.6cm}p{3.2cm}p{4.4cm}@{}}
\toprule
\textbf{Operation} & \textbf{Type} & \textbf{Notes} \\
\midrule
\texttt{select}, \texttt{filter}, \texttt{withColumn} & Narrow & No shuffle; stays in partition \\
\texttt{groupBy().agg()} & Wide & Shuffle by group key \\
\texttt{join()} & Wide & Shuffle unless broadcast or co-located \\
\texttt{broadcast(dim)} & Hint & Avoids the shuffle for small dims \\
\texttt{distinct()}, \texttt{dropDuplicates()} & Wide & State/shuffle on the full key \\
\texttt{repartition(n)} / \texttt{coalesce(n)} & Wide & Reshape partitions before writes \\
\texttt{cache()} & Eager marker & Reuse across multiple actions only \\
\bottomrule
\end{tabular}
\end{center}

Key configuration: \texttt{spark.sql.shuffle.partitions} (default 200---tune to
data volume), \texttt{spark.sql.adaptive.enabled} (skew and coalesce handling),
and Photon on Databricks for vectorized scans (Chapter 11).

\section{Delta Lake Operations}
\index{Delta Lake!operations}

\begin{lstlisting}[language=SQL,caption={Delta maintenance and time travel (Spark SQL).},label={lst:app-delta}]
MERGE INTO silver.orders AS tgt
USING staging.batch AS src ON tgt.order_id = src.order_id
WHEN MATCHED AND src.updated_at > tgt.updated_at THEN UPDATE SET *
WHEN NOT MATCHED THEN INSERT *;

SELECT * FROM silver.orders VERSION AS OF 12;
SELECT * FROM silver.orders TIMESTAMP AS OF '2026-08-01T00:00:00';

DESCRIBE HISTORY silver.orders;
DESCRIBE DETAIL  silver.orders;   -- file counts, sizes, partitioning

OPTIMIZE silver.orders ZORDER BY (customer_id);
VACUUM silver.orders RETAIN 168 HOURS;   -- 7 days; check time-travel consumers first
\end{lstlisting}

\begin{pitfall}
\textbf{VACUUM under the retention horizon of a time-travel consumer.} Files
removed by VACUUM break \texttt{VERSION AS OF} queries older than the retention.
Coordinate the horizon with audit and reprocessing requirements (Chapters 6, 23).
\end{pitfall}

\section{Structured Streaming}
\index{Structured Streaming!reference}

\begin{lstlisting}[language=Python,caption={The Chapter 15 streaming skeleton, condensed.},label={lst:app-stream}]
stream = (spark.readStream.format("kafka").options(**kafka_options).load()
    .select(F.from_json(F.col("value").cast("string"), schema).alias("e"))
    .select("e.*")
    .withColumn("event_time", F.to_timestamp("ts")))

agg = (stream.withWatermark("event_time", "10 minutes")
       .groupBy(F.window("event_time", "5 minutes"), "deviceId")
       .agg(F.avg("tempC").alias("avg_temp")))

(agg.writeStream.format("delta").outputMode("append")
    .option("checkpointLocation", "abfss://checkpoints@<adls>/iot/v1")
    .trigger(processingTime="1 minute")     # or availableNow=True for batch-style
    .toTable("iot_agg"))
\end{lstlisting}

Remember: version checkpoint paths per deployment; never restart changed logic
from a stale checkpoint.

\section{Synapse T-SQL Essentials}
\index{Synapse!T-SQL reference}

\begin{lstlisting}[language=SQL,caption={Dedicated and serverless pool statements used in Chapters 5--8.},label={lst:app-tsql}]
-- bulk load (dedicated pool)
COPY INTO stg_sales
FROM 'https://<adls>.dfs.core.windows.net/raw/sales/*.parquet'
WITH (FILE_TYPE='PARQUET', CREDENTIAL=(IDENTITY='Managed Identity'));

-- distributed table creation
CREATE TABLE fact_sales
WITH (DISTRIBUTION = HASH(customer_id), CLUSTERED COLUMNSTORE INDEX)
AS SELECT * FROM stg_sales;

-- serverless: query the lake in place
SELECT * FROM OPENROWSET(
    BULK 'https://<adls>.dfs.core.windows.net/gold/fact_sales/',
    FORMAT = 'DELTA') AS f
WHERE f.filepath(1) = '2026';

-- serverless: materialize back to the lake
CREATE EXTERNAL TABLE curated.orders_clean
WITH (LOCATION='orders_clean/', DATA_SOURCE=ds_curated,
      FILE_FORMAT=parquet_ff)
AS SELECT ... ;

-- in-database scoring (dedicated pool)
SELECT * FROM PREDICT(MODEL = @model, DATA = dbo.features AS d,
                      RUNTIME = ONNX) WITH (score REAL);
\end{lstlisting}

\section{Useful DMVs}
\index{DMV}
\begin{center}
\footnotesize
\begin{tabular}{@{}p{6.4cm}p{6.8cm}@{}}
\toprule
\textbf{DMV} & \textbf{Answers} \\
\midrule
\texttt{sys.dm\_pdw\_exec\_requests} & What ran, how long, which label \\
\texttt{sys.dm\_pdw\_request\_steps} & Where data moved (shuffle steps) \\
\texttt{sys.pdw\_nodes\_column\_store\_row\_groups} & Rowgroup health per table \\
\texttt{sys.dm\_pdw\_wlm\_waiting\_requests} & Concurrency queue depth \\
\texttt{sys.workload\_management\_workload\_classifiers} & Classifier routing rules \\
\bottomrule
\end{tabular}
\end{center}

\section{Cross-References}

Chapters 5--7 (warehouse and lakehouse SQL), 9 and 12 (Spark and Delta), 15
(streaming), and 17 (PREDICT). Primary documentation: Spark \cite{apache2025spark}, Structured
Streaming \cite{apache2025structuredStreaming}, Delta Lake
\cite{armbrust2020delta,deltaio2025delta}, and Synapse best practices
\cite{microsoft2025synapseDedicatedBest,microsoft2025synapseServerless}.
